\section{Signals background: how we get complex numbers out of an antenna}
\label{sec:iq}

The math section ahead manipulates ``complex baseband samples'' as if that
were the obvious thing a radio produces. It is not obvious. This section
motivates it from scratch: what the antenna really measures, why we describe
it with complex numbers, what I/Q means, and precise working definitions of
the terms (narrowband, broadband, window, coherent) used everywhere else.

\subsection{What the antenna actually measures}
\label{sec:antenna}

An antenna is a wire; the only thing it can give the receiver is a
\emph{real} voltage that wiggles in time, $v(t)$. For a 2.412\,GHz emitter
that voltage oscillates 2.4 \emph{billion} times per second. Almost none of
those wiggles are interesting individually --- the information (and, for us,
the geometry) rides on how the wiggle's \emph{strength} and \emph{timing}
drift over microseconds, millions of times slower than the wiggle itself.

Any such signal can be written as
\[
v(t) \;=\; A(t)\,\cos\!\big(2\pi f t + \psi(t)\big),
\]
where $f$ is the \textbf{carrier frequency} (the fast wiggle),
$A(t) \ge 0$ is the \textbf{amplitude} or \textbf{envelope} (how strong the
wiggle currently is), and $\psi(t)$ is the \textbf{phase} (how the wiggle is
currently shifted relative to a perfect clock at frequency $f$). $A$ and
$\psi$ change slowly compared to the carrier; they are ``the signal''. The
pair $(A, \psi)$ is what we actually want, sampled in time.

\subsection{Why complex numbers: one arrow instead of two numbers}

A pair (strength, angle) is naturally a single \emph{arrow} in the plane ---
a complex number:
\[
x(t) \;=\; A(t)\, e^{j\psi(t)}
\qquad\text{(length $=$ amplitude, direction $=$ phase).}
\]
This is not just notation. The two operations radio processing does
constantly become trivial arithmetic on arrows:
\begin{itemize}
  \item \textbf{Delaying} a narrowband signal by $\tau$ rotates its arrow by
        a fixed angle: $x(t-\tau) \approx x(t)\,e^{-j2\pi f \tau}$
        (the envelope barely moves, the carrier phase shifts). Delay ---
        which is geometry --- becomes multiplication.
  \item \textbf{Comparing} two channels becomes one multiplication:
        $x_1 x_0^{*}$ has length $A_1 A_0$ and angle $\psi_1 - \psi_0$. The
        entire two-antenna measurement of Section~\ref{sec:math} is this one
        product.
\end{itemize}
With real cosines both operations are trigonometric-identity soup; with
arrows they are one-liners. That is the whole reason radios use complex
numbers.

\subsection{I/Q: how the hardware extracts the arrow}
\label{sec:iq-demod}

\paragraph{The cast.}
Recall from Section~\ref{sec:antenna} what $v(t)$ is: the \emph{real} voltage
on the antenna wire,
\[
v(t) = A\cos(2\pi f t + \psi),
\]
where $f$ is the carrier ($\sim$2.4 billion wiggles per second), $A$ is how
strong the wave currently is, and $\psi$ is its timing shift relative to a
perfect clock at frequency $f$. ($A$ and $\psi$ do change with time, but a
million times slower than the carrier --- over the few nanoseconds this
derivation cares about they are constants.) The receiver wants $A$ and $\psi$
but can only measure $v(t)$, in which they are tangled up inside the cosine.
Three hardware pieces untangle them:
\begin{itemize}
  \item the \textbf{local oscillator (LO)}: the radio's own clean internal
        tone at the carrier frequency --- it can generate $\cos(2\pi f t)$ and
        (by delaying a quarter cycle) $\sin(2\pi f t)$;
  \item a \textbf{mixer}: a circuit that literally multiplies two voltages;
  \item a \textbf{low-pass filter}: a circuit that averages away anything
        wiggling faster than some cutoff, keeping only slow variations.
\end{itemize}

\paragraph{The one identity that does all the work.}
Multiplying two cosines produces the sum and difference of their angles ---
the product-to-sum identity from trigonometry:
\[
2\cos(a)\cos(b) = \cos(a-b) + \cos(a+b),
\qquad
2\cos(a)\sin(b) = \sin(a+b) - \sin(a-b).
\]
The trick of the whole receiver: choose $b$ = the LO's angle $2\pi f t$, and
$a$ = the signal's angle $2\pi f t + \psi$. Then the \emph{difference}
$a - b = \psi$ --- the fast carrier subtracts out, leaving exactly the
quantity we want --- while the \emph{sum} $a + b = 4\pi f t + \psi$ wiggles at
twice the carrier and can be filtered away.

\paragraph{I branch, line by line.}
Multiply the antenna voltage by twice the LO cosine and apply the first
identity with $a = 2\pi f t + \psi$, $b = 2\pi f t$:
\begin{align*}
v(t)\cdot 2\cos(2\pi f t)
  &= 2A\cos(2\pi f t + \psi)\cos(2\pi f t)\\
  &= A\cos(\underbrace{\psi}_{a-b})
     \;+\; \underbrace{A\cos(4\pi f t + \psi)}_{a+b:\ \text{wiggles at }2f}
  \;\xrightarrow{\text{low-pass}}\; I = A\cos\psi .
\end{align*}
The second term oscillates at $2f \approx 4.8$\,GHz --- averaging over even a
few carrier cycles kills it; the first term is constant (on carrier
timescales) and sails through the filter.

\paragraph{Q branch, line by line.}
Same signal, multiplied by the quarter-cycle-shifted copy of the LO, using
the second identity:
\begin{align*}
v(t)\cdot \big(-2\sin(2\pi f t)\big)
  &= -2A\cos(2\pi f t + \psi)\sin(2\pi f t)\\
  &= A\sin(\underbrace{\psi}_{a-b})
     \;-\; \underbrace{A\sin(4\pi f t + \psi)}_{\text{at }2f}
  \;\xrightarrow{\text{low-pass}}\; Q = A\sin\psi .
\end{align*}
(The minus sign on the LO copy is chosen so the surviving term comes out as
$+A\sin\psi$.)

\paragraph{Why two branches are necessary.}
One branch alone gives $I = A\cos\psi$ --- a single number that hopelessly
mixes strength and timing: $I = 0.5$ could be a strong signal at
$\psi = 60^\circ$ or a weak one at $\psi = 0^\circ$, and $\psi$ and $-\psi$
are indistinguishable. The arrow lives in a plane; you need its shadow along
\emph{two} perpendicular directions to pin it down. $I$ and $Q$ are exactly
those two shadows, and together they reconstruct it uniquely:
\[
x = I + jQ = A\cos\psi + jA\sin\psi = A e^{j\psi},
\qquad
A = \sqrt{I^2 + Q^2}, \quad \psi = \operatorname{atan2}(Q, I).
\]
Worked example: $A = 1$, $\psi = 60^\circ$ gives $I = 0.5$, $Q = 0.866$ ---
the point $(0.5, 0.866)$, an arrow of length 1 at $60^\circ$.

\textbf{I} is the \emph{in-phase} component (relative to the LO cosine),
\textbf{Q} the \emph{quadrature} (90$^\circ$-shifted) component. The SDR
digitizes both at the \textbf{sample rate} $f_s$ (up to 30\,MHz here),
producing the complex \textbf{baseband} stream $x[n]$ used everywhere in this
document. ``Baseband'' means the carrier has been divided out: a pure
unmodulated carrier at exactly the LO frequency becomes a \emph{constant}
arrow.

\paragraph{When the LO is not exactly the carrier.}
Everything above assumed the LO frequency equals $f$. In reality the radio is
tuned to $f_{\mathrm{LO}} = f - f_{\mathrm{IF}}$ (deliberately, plus crystal
error), and redoing the derivation with $b = 2\pi f_{\mathrm{LO}} t$ gives
$a - b = 2\pi f_{\mathrm{IF}} t + \psi$: the surviving slow term is not a
constant but a slow rotation at the \emph{intermediate frequency}
$f_{\mathrm{IF}} = f - f_{\mathrm{LO}}$. The arrow spins at $f_{\mathrm{IF}}$
turns per second. This is why IF placement matters so much for data quality
(the dataset-quality report's tone-at-DC analysis): $f_{\mathrm{IF}} \approx 0$
parks the spinning arrow on top of the receiver's DC machinery, while a
healthy offset keeps it clear --- and since both antennas share one LO, the
spin is common to both channels and cancels in the phase difference.

Two hardware facts matter enormously for SPF:
\begin{itemize}
  \item \textbf{Both RX channels of the AD9361 share one LO.} The arbitrary
        LO phase appears identically in $\psi_0$ and $\psi_1$ and cancels in
        $x_1 x_0^{*}$. This is why the \emph{inter-channel phase difference}
        is physically meaningful and stable, even across signal bursts, while
        the \emph{absolute} phase of either channel alone is meaningless.
  \item \textbf{The LO is never exactly the emitter's carrier.} A small
        frequency offset makes the arrow rotate slowly (``spin''); it is
        common to both channels and cancels in the difference, but it is why
        single-channel phase drifts and must be detrended in some analyses.
\end{itemize}

\paragraph{Remark: I/Q as value and derivative (phase space).}
There is a physicist's way to see all of this at once. At any instant, the
waveform's \emph{value} is $v = A\cos\theta$ and its \emph{scaled
derivative} is $-v'/(2\pi f) = A\sin\theta$ (differentiating a cosine
yields the sine; the carrier dominates the slope). So (value, derivative) are
the same two numbers as (I, Q), just measured in the lab frame at the total
angle $\theta = 2\pi f t + \psi$; the LO's rotating frame subtracts the
known carrier ramp and freezes the arrow at $\psi$. Deeper still: a sinusoid
is a harmonic oscillator, whose state space is exactly two-dimensional ---
(position, velocity), $(A, \psi)$, and $(I, Q)$ are three coordinates for the
same state, which is why two branches are necessary and sufficient (and why
$I$ and $Q$ are called \emph{quadratures}, the oscillator's $x$ and $p$).
The hardware mixes rather than differentiates for practical reasons only: a
differentiator at 2.4\,GHz would need sub-cycle timing, amplifies noise
$\propto$ frequency, and reads one instant --- while mixer + low-pass
measures the same two numbers averaged over thousands of carrier cycles.

\subsection{Term definitions used throughout}
\label{sec:terms}

\begin{description}
  \item[Carrier / baseband.] The carrier is the fast oscillation at $f$;
        baseband is the slow complex signal $x(t)$ left after the carrier is
        mixed out.
  \item[Bandwidth $B$.] How fast the arrow moves: the width of the frequency
        band occupied by $A(t)e^{j\psi(t)}$. A pure tone has $B \to 0$; a
        20\,MHz Wi-Fi packet has $B = 20$\,MHz.
  \item[Narrowband (for our purposes).] Two conditions, both comfortably true
        in SPF: (i) $B \ll f$ (the signal is a thin slice around the
        carrier), and (ii) --- the one that matters for arrays ---
        $B\,\tau \ll 1$, where $\tau \le d/c$ is the inter-antenna delay: the
        envelope does not change in the $\sim$0.2\,ns it takes the wave to
        cross the array, so both antennas see \emph{the same} $A$ and $\psi$
        and differ only by the geometric rotation $e^{-j2\pi f\tau}$. Even
        Wi-Fi ($B\tau \approx 4\times10^{-3}$) passes easily.
  \item[Broadband / wideband.] $B\tau$ not small, or $B/f$ large: different
        frequency components inside the signal have different wavelengths,
        hence \emph{different} inter-antenna phase shifts --- a single
        $\varphi$ no longer summarizes the geometry and processing must split
        the band (e.g.\ per-subcarrier). SPF's beacons do not operate here,
        but strong wideband interferers do.
  \item[Window.] A contiguous run of $N$ samples treated as one measurement
        unit; SPF computes phase/beamformer statistics per window, many
        windows per snapshot.
  \item[Snapshot / session.] A snapshot is one capture buffer with one
        ground-truth pose; a session (dataset) is the full recorded
        trajectory of snapshots.
  \item[Coherent vs.\ incoherent combining.] Coherent: add the complex
        arrows, \emph{then} take magnitude/angle --- signals aligned in phase
        reinforce ($N$ arrows of unit length aligned give length $N$; random
        phases give $\sim\sqrt{N}$), so SNR grows with window length.
        Incoherent: add magnitudes or powers, discarding phase --- robust
        when phase is unstable, but no phase gain. The coherent sum
        $\sum x_1 x_0^{*}$ of Section~\ref{sec:estimation} is the coherent
        combine of the \emph{difference} arrows.
  \item[SNR.] Signal-to-noise power ratio within the processed band; window
        statistics turn per-sample SNR into per-window estimate quality.
\end{description}

With this vocabulary, the math section can be read literally: every
``complex baseband sample'' is an I/Q arrow, every phase difference is an
angle between two arrows sharing one clock, and every ``narrowband''
invocation is the statement $B\tau \ll 1$ justified above.
